\section{Hardening Democracy Against Populism} \label{sec:hardening_democracy}

The struggle against populism must take place on several levels.
It can only truly be defeated by depriving it of its substantive basis.

Apparently, however, this is not so easy, because democracy, too, produces elites and, much worse, a disconnection of the elites from the everyday reality of the non-elites.
The existence of elites as such is rather a good thing: none of us wants to be led by people who are dumber than we are.
Actually, we do not want to be led by people who are just as clever as we are either.
Instead, we want to be led by the best among us.

It becomes problematic when ever larger parts of the population gain the impression that
\begin{enumerate}[label=\alph*)]
    \item the elites live in their own world that no longer has anything to do with reality, or
    \item only pursue their own interests, which run counter to the interests of the \enquote{ordinary citizen}, or
    \item cannot read the times correctly, or
    \item simply have no solution either, or
    \item the elites are not really cleverer than they are, but merely privileged.
\end{enumerate}

The only real way to counter this is by once again bringing people from the bottom to the top.
Unlike monarchy, democracy does allow this without revolution.

However, it does not necessarily make it easy to rise from the bottom to the top either.
With wealthy, educated parents in the right party, one has considerably better chances of advancement than a child from a background far removed from education and politics.

Nevertheless, it is possible for people who know real life to make it to the top.
But it may take a while until someone once again manages to reach a leadership position who is truly a good leader, and not just someone who is a good talker or has the necessary connections.

It is also difficult, though not impossible, to found new parties and let them grow to such a size that they are genuine electable alternatives to the existing ones.

Because democracy permits bloodless upheavals but does not make them easy either, a temporary weakness in leadership, in terms of both personnel and substance, can arise.
This is then easily exploited by populists, people who claim to have simple solutions but in reality make things even worse.

Populists come in two varieties:
\begin{enumerate}[label=\alph*)]
    \item Democratic populists. People with dangerously stupid ideas who nevertheless do not question democracy.
    \item Undemocratic populists. People with dangerously stupid ideas who want to turn democracy into an autocracy.
\end{enumerate}

The way from a) to b) is not far.
When you realize that your stupid ideas do not get through, you can quickly blame democracy and move up from a) to b), as can be seen clearly in the example of Donald Trump's first and second terms.

Whether, and if so how, democracy could be improved further so that positive upheavals would become easier without making negative upheavals easier at the same time, I do not know.
But one can try to make democracy resilient against temporary phases of \enquote{rule of fools}, so as to give the people a little more time to bring forth better leaders and, if necessary, to correct a mistake.

What matters is that the phases are temporary.
The assumption is that the people, when sufficiently desperate about those who have governed so far, can quickly vote a populist party into office, but, like America under Trump's first term, soon realize that they have made a mistake.\footnote{In my assessment, Trump's later re-election is not a correction of the correction, but came about rather through a combination of forgetfulness and despair over the lack of electoral alternatives.}

The democratic structures must then be so stable that one legislative term is not enough to rebuild them, so that in the end there are still free elections at which the party can be voted out again.
It would also be good if the very worst decisions of a populist-dominated Bundestag and a populist Chancellor could still be blocked from outside the Bundestag.

If populist rule extends over several legislative terms, they will probably manage to rebuild democracy according to their ideas over time after all.
It is in the nature of democracy that the people get what they want.
So if they no longer want democracy, they get none.

In a democracy, the people cannot be protected from themselves in the long run.
That is only possible in autocracies.
But they can be protected from the consequences of a lapse.
